% ---- pg-262 （书页 250） ----
对已婚男女各 30 名($n_1=30, n_2=30$ 的两个样本)进行了调查, 得出样本均值向量和融合的样本协方差阵如下:
\[
\bar{\mathbf{x}}_1=\underset{(\text{男性})}{\begin{bmatrix}6.833\\7.033\\3.967\\4.700\end{bmatrix}}, \qquad
\bar{\mathbf{x}}_2=\underset{(\text{女性})}{\begin{bmatrix}6.633\\7.000\\4.000\\4.533\end{bmatrix}}
\]
\[
\mathbf{S}_p=\begin{bmatrix}
0.606&0.262&0.066&0.161\\
0.262&0.637&0.173&0.143\\
0.066&0.173&0.810&0.029\\
0.161&0.143&0.029&0.306
\end{bmatrix}
\]
样本均值向量作为样本轮廓如\hyperref[fig:6.5]{图 6.5} 所示.

\begin{mvfigure}
\begin{tikzpicture}[scale=0.78]
% 坐标轴
\draw (0,1.6)--(0,7.7);
\draw[->] (0,1.6)--(5.1,1.6) node[right,font=\small]{变量};
% y 轴断口
\draw (0,2.35)--(-0.13,2.2)--(0.13,2.05)--(-0.13,1.9)--(0,1.75);
% y 轴刻度
\foreach \y in {2,4,6} \draw (-0.08,\y)--(0.08,\y) node[left,font=\small]{\y};
\node[rotate=90,font=\small] at (-0.45,4.6) {样本均值$\bar{x}_{i}$};
% x 轴刻度
\foreach \x in {1,2,3,4} \draw (\x,1.6)--(\x,1.48) node[below,font=\small]{\x};
% 男性折线
\draw[line width=0.9pt] (1,6.55)--(2,6.95)--(3,3.75)--(4,4.65);
\foreach \p in {(1,6.55),(2,6.95),(3,3.75),(4,4.65)} \node[font=\small] at \p {$\times$};
% 女性折线
\draw[line width=0.9pt] (1,6.45)--(2,6.85)--(3,3.95)--(4,4.58);
\foreach \p in {(1,6.45),(2,6.85),(3,3.95),(4,4.58)} \draw \p circle (2.2pt);
% 图例
\draw (1.3,3.35)--(1.9,3.35); \node[font=\small] at (1.6,3.35) {$\times$};
\node[right,font=\small] at (1.9,3.35) {男性};
\draw (1.3,2.95)--(1.9,2.95); \node[font=\small] at (1.6,2.95) {$\circ$};
\node[right,font=\small] at (1.9,2.95) {女性};
\end{tikzpicture}
\phantomsection\label{fig:6.5}\mvcap{图 6.5\quad 对爱情—婚姻问题的回答所构成的样本轮廓}
\end{mvfigure}

尽管回答均取整数值, 它们虽然不是正态随机变量, 但由于样本容量足够大, 我们仍可采用正态理论的方法来处理. 为检验轮廓的平行性(即 $H_{01}: \mathbf{C}\boldsymbol{\mu}_1=\mathbf{C}\boldsymbol{\mu}_2$), 我们算出
\[
\mathbf{C}\mathbf{S}_p\mathbf{C}'=\begin{bmatrix}-1&1&0&0\\0&-1&1&0\\0&0&-1&1\end{bmatrix}\cdot\mathbf{S}_p\cdot\begin{bmatrix}-1&0&0\\1&-1&0\\0&1&-1\\0&0&1\end{bmatrix}
=\begin{bmatrix}0.719&-0.268&-0.125\\-0.268&1.101&-0.751\\-0.125&-0.751&1.058\end{bmatrix}
\]
和
\[
\mathbf{C}(\bar{\mathbf{x}}_1-\bar{\mathbf{x}}_2)=\begin{bmatrix}-1&1&0&0\\0&-1&1&0\\0&0&-1&1\end{bmatrix}\begin{bmatrix}0.200\\0.033\\-0.033\\0.167\end{bmatrix}=\begin{bmatrix}-0.167\\-0.066\\0.200\end{bmatrix}
\]
因而
